\scanpage{99}
kif-kif, car les composantes connexes ne changent pas en ajoutant ou retranchant des branches orphelines), on trouve la relation la plus générale entre $G$ et un $G'_0$, du type explicité dans la définition des ``morphismes de normalisation de m.graphes''.

On définit de façon évidente les composés de morphismes de normalisation. Le composé de morphismes de normalisation strict est strict. Les m.graphes avec morphismes de normalisation strict font une catégorie, dont les iso. sont les iso. ordinaires de m.graphes. Les morphismes de cette catégorie sont des monomorphismes. Quand on se borne aux morphismes \inserted{de norm.} stricts, on voit que (comme pour les morphismes de générisation) l'ens des sous-m.graphes stricts d'un m.graphe \inserted{$G$} est un ens. ordonné \deleted{\unsure{iso}}\inserted{anti}morphe à l'ens des parties de $A(G)$, en associant à tte partie $C\subset A(G$
\margin{98}
